
\chapter{Жизненный цикл транзакции}\label{ch:tx-lifecycle}
Покупатель прикладывает карту в~кофейне, терминал показывает
\enquote{одобрено}. Деньги продавец увидит на счёте позже, иногда через
несколько суток. Карточная транзакция проходит три фазы~--- авторизацию,
клиринг, расчёт~--- и~сумма в~них не обязана совпадать. На стыках фаз возникают
\emph{три разрыва}, на которых построена ежедневная сверка
(см.~раздел~\ref{sec:recon-gaps}).

\section{Авторизация}\label{sec:tx-authorization}

Авторизация~--- единственная фаза карточного платежа, которая идёт в~реальном
времени и~видна покупателю. Для покупателя контакт с~платёжной системой
заканчивается на ответе терминала; клиринг, расчёт и~возможный диспут проходят
позже без его участия. В~эти несколько секунд фиксируются реквизиты, по~которым
затем пойдут клиринг и~спор.

\subsection*{Путь авторизационного запроса}

Когда держатель прикладывает карту к~терминалу, вставляет её или проводит
магнитной полосой, терминал считывает данные карты и~формирует авторизационный
запрос. Это структурированное сообщение в~формате ISO~8583; подробно о~формате
см.\ в~главе~\ref{ch:iso8583}. Внутри собраны как минимум PAN в~DE~2 (Data
Element, элемент данных) и~данные Track~2 в~DE~35 или их чиповый эквивалент
в~DE~55, сумма и~валюта операции в~DE~4 и~DE~49, идентификаторы терминала и~ТСП
в~DE~41 и~DE~42, категория ТСП~--- MCC (Merchant Category Code)~--- в~DE~18,
а~также тип операции, который задают MTI (Message Type Indicator, тип
сообщения) и~processing code в~DE~3.

Сообщение уходит эквайеру~--- банку или процессору, обслуживающему этот
терминал. Эквайер выполняет первичные проверки: валидность терминала,
соответствие MCC, базовые правила антифрода на своей стороне. Затем он
маршрутизирует запрос в~карточную сеть~\cite{visa-core-rules}. Сеть по IIN/BIN
определяет эмитента или его процессор и~передаёт запрос адресату.

Весь путь от~терминала к~эквайеру, в~сеть, эмитенту и~обратно должен уложиться
в~жёсткие сроки. У Visa публичные правила задают верхние пределы ожидания для
авторизаций с~POS-терминалов (POS, Point-of-Sale, точка продажи) порядка
нескольких секунд, с~более коротким пределом в~Европе. Таймауты для снятия
наличных в~банкомате (ATM, Automated Teller Machine) задаются отдельной строкой
в~таблице \enquote{Maximum Time Limits for Authorization Request Response}
актуальной редакции Visa Core Rules~\cite{visa-core-rules}. Если ответ не
пришёл вовремя, за эмитента может ответить сама Visa через механизм Stand-In
Processing (STIP). STIP одобряет операции с~низким профилем риска: сумма
в~пределах нормы для карты и~ТСП, знакомая география, отсутствие признаков
фрода. Неясная операция чаще отклоняется, потому что в~STIP риск переходит
к~сети; поэтому механизм работает только для низкорисковой части потока. У Mastercard
пределы ожидания и~показатели доступности задаются в~спецификациях сообщений
и~контролируются через нормативы доступности
авторизации~\cite{visa-core-rules,mc-tpr}.

\subsection*{Решение эмитента}

Получив авторизационный запрос, эмитент или его процессор решает одобрить или
отклонить операцию по нескольким группам проверок, которые выполняются за доли
секунды:

\begin{itemize}
  \item \textbf{проверка карты}~--- существует ли карта с~таким PAN? Не
    заблокирована ли она? Не истёк ли срок действия? Совпадает ли CVV или
    криптограмма?
  \item \textbf{проверка баланса}~--- достаточно ли средств на счёте дебетовой
    карты или доступного кредитного лимита у~кредитной?
  \item \textbf{лимиты и~ограничения}~--- не превышен ли дневной лимит
    на~операции? Не заблокирована ли карта для этого типа операций,
    например запрет на~интернет-покупки? Не ограничена ли география
    использования?
  \item \textbf{антифрод}~--- соответствует ли операция типичному поведению
    держателя? Нет ли признаков компрометации карты? Не входит ли ТСП
    в~списки повышенного риска?
  \item \textbf{аутентификация держателя}~--- для чиповых транзакций верификация
    криптограммы ARQC (Authorization Request Cryptogram,
    см.~гл.~\ref{ch:emv}); для CNP (Card-Not-Present, без физического
    предъявления карты)~--- результат 3-D Secure, если он был; для
    магнитной полосы~--- CVV.
\end{itemize}

По результатам этих проверок эмитент формирует ответ с~кодом ответа в~DE~39.
Значение \texttt{00} означает, что операция одобрена; \texttt{05}~--- общий
отказ без детализации; \texttt{51}~--- недостаточно средств; \texttt{14}~---
неверный номер карты. Полный список кодов зависит от
ISO~8583-профиля и~правил конкретной сети; базовая модель стандартизована, но
сети дополняют её собственными значениями
и~уточнениями~\cite{iso-8583-2023,visa-tadg,mc-rules}. Классификация кодов по
следующему действию (жёсткий, мягкий финансовый, аутентификационный,
технический, непрозрачный) и~правила повторов разобраны
в~разделе~\ref{sec:decline-response-codes}.

Если транзакция одобрена, эмитент присваивает ей код авторизации в~DE~38~---
шестисимвольный идентификатор, привязывающий одобрение к~конкретному запросу.
На этапе клиринга именно по нему сопоставят авторизацию с~финансовой записью.
Одновременно эмитент блокирует запрошенную сумму на счёте держателя~--- холд
(\textit{hold}). Деньги ещё не списаны, но уже недоступны для других операций.

\subsection*{Срок жизни холда}

Правила карточных сетей задают окно между авторизацией и~моментом, когда
транзакцию нужно либо представить в~клиринг, либо снять отменой
(\textit{reversal}). У Visa для большинства транзакций с~физическим
предъявлением карты это окно составляет 5~календарных дней; для операций без
предъявления~--- 10~дней; для оценочной авторизации в~гостиницах, круизах
и~аренде автомобилей~--- до~30~дней. У Mastercard клиринговое окно зависит от
типа авторизации (final, undefined, preauthorization) и~находится в~диапазоне
7--30~календарных дней~--- пункт~3.15.1 \enquote{Transaction Presentment Time
Frames} Mastercard Transaction Processing Rules. Эти окна и~задают, как долго
эмитент держит неиспользованный резерв; точный момент снятия холда зависит ещё
и~от внутренних правил эмитента~\cite{visa-core-rules,mc-tpr}.

Числа 5,~7 и~30~дней отражают минимальный операционный цикл: время доставки или
оказания услуги, передачу клирингового файла от эквайера в~сеть и~буфер на
задержки между участниками. Длинные сроки для гостиниц и~аренды автомобилей
соответствуют реальной дельте между авторизацией при заселении
(\textit{check-in}) и~финальной суммой при выезде. Эти сроки исторически
сложились как компромисс сетей; их изменение требует согласования со~всеми
участниками расчёта.

В~системе эмитента холд~--- это запись в~таблице, синхронизация
которой с~доступным остатком держится на корректной последовательности коммитов
авторизации и~отката. Любой сбой в~этой цепочке оставляет клиенту видимое
расхождение даже без участия продавца.

\practicenote{%
  Если клиент жалуется, что \enquote{деньги списали, но покупка
  не прошла}, чаще всего речь идёт о~холде. Авторизация была
  одобрена и~средства заблокированы, но продавец
  не завершил операцию: например, отменил заказ
  на своей стороне, не отправив отмену. Деньги вернутся
  автоматически по истечении холда, а~в~мобильном банке клиент до этого
  видит уменьшенный доступный баланс.

  Сложнее, когда расхождение возникает внутри хоста эмитента. Частичный отказ
  обработки рассинхронизирует холд и~счёт: запрос одобрен, обновление таблицы
  холдов не сохранилось~--- резерв остаётся без нормального срока снятия при
  видимом клиенту отказе. Обратный случай: холд снят при откате незафиксированной
  транзакции, и~клиент тратит средства, которые ещё \texttt{pending}. Защита~---
  отдельный журнал холдов с~меткой времени авторизации и~признаком сверки
  с~клирингом.
}

\section{Клиринг}\label{sec:tx-clearing}

\highlight{Авторизация~--- обещание заплатить, клиринг~--- предъявление счёта,
формальное требование эквайера к~эмитенту исполнить это обещание}. Клиринг
инициирует эквайер: он формирует у~себя пакет финансовых записей и~отправляет
его в~сеть, а~эмитент получает пакет как адресат. Клиринг идёт пакетами
(batches), без секундного ожидания ответа; покупатель в~этом цикле уже
не~участвует.

\subsection*{Формирование клирингового файла}

В~конце операционного дня~--- а~у~некоторых процессоров и~несколько раз
за~день~--- эквайер формирует клиринговый файл (\textit{clearing file}): пакет
записей обо всех транзакциях, прошедших авторизацию за период. Каждая запись
содержит реквизиты транзакции: PAN, сумму, которая может отличаться от
авторизованной, валюту, дату и~время, идентификатор ТСП, MCC, код авторизации
и~другие данные, нужные для расчёта~\cite{mc-rules}. Отправку такого файла
с~итоговыми суммами называют финансовым представлением (\textit{financial
presentment}) или захватом (\textit{capture}); в~этот момент эквайер фиксирует
финальную сумму обязательства~\cite{visa-tadg,mc-tpr}.

Формат клирингового файла зависит от карточной сети. Visa использует BASE~II,
в~котором каждый тип записи обозначается Transaction Code (TC); Mastercard~---
формат IPM (Integrated Product Messages) с~MTI~1240 для~first presentment; НСПК
для~внутрироссийских операций по~картам \enquote{Мир} и~по~картам международных
сетей~--- собственный профиль ISO~8583 с~расширениями НСПК. Состав полей и~коды
причин относятся к~закрытой части документации, доставка идёт пакетным файлом
через систему электронного документооборота НСПК~\cite{nspk-mir-rules}. Логика
у~всех одна: сеть получает набор записей с~реквизитами транзакций, по которым
нужно провести расчёт~\cite{visa-core-rules,mc-rules}.

\subsection*{Сопоставление с~авторизацией}

Получив клиринговый файл, сеть сопоставляет каждую запись с~соответствующей
авторизацией по~комбинации PAN, суммы, кода авторизации и~даты. Сумма
в~клиринге не~обязательно равна авторизованной~--- возможны три исхода:

\begin{itemize}
  \item \textbf{auth = clearing} (штатный сценарий, большинство покупок)~---
    покупка по~POS в~рознице: терминал формирует авторизационный запрос
    на~финальную сумму, ровно она и~приходит в~клиринг.
  \item \textbf{auth > clearing}~--- авторизация была больше итогового чека.
    Классические случаи~--- предавторизация отеля или АЗС
    на~ориентировочную сумму, фактический счёт меньше. Эквайер выпускает
    частичную отмену (\textit{partial reversal}) до~клиринга, либо
    клиринговая запись приходит сразу на~финальную сумму; разница
    освобождается у~держателя при снятии холда.
  \item \textbf{auth < clearing}~--- клиринговая сумма выше авторизованной.
    Типичный пример~--- ресторан, где авторизация проходит на~сумму заказа,
    а~в~клиринг попадает сумма с~чаевыми. Правила сетей допускают такое
    превышение для~чаевых и~категорий travel\&entertainment,
    но~ограничивают его процентом или фиксированным лимитом; точные пределы
    зависят от~MCC, программы и~текущей редакции
    правил~\cite{visa-core-rules,mc-rules}.
\end{itemize}

Если представление в~клиринг приходит позже применимого окна или данные в~нём
не совпадают с~исходной авторизацией, сеть всё равно может доставить это
представление эмитенту. Но для эквайера это позднее представление (\textit{late
presentment}): риск диспута повышается, и~эквайер теряет часть защитного срока
по спору~\cite{visa-core-rules,mc-tpr,mc-chargeback-guide}.

Более опасный сбой: клиринговый файл теряется или не принимается сетью из-за
ошибки формата. Эквайер узнаёт об этом, когда сеть возвращает отчёт с~перечнем
отклонённых записей~--- \textit{rejected transactions}. Эквайер
исправляет записи и~повторно отправляет файл в~следующем пакете клиринга.
Правила сетей ограничивают срок повторной подачи. У Visa \textit{financial
presentment} должен попасть в~сеть в~пределах авторизационного окна (5--30~дней
в~зависимости от MCC), иначе авторизация аннулируется~\cite{visa-core-rules}. У
Mastercard клиринг должен прийти в~пределах 7--30~календарных
дней~\cite{mc-tpr}. Если эквайер не уложился, эмитент автоматически снимает
холд, и~деньги возвращаются в~доступный баланс держателя. Для продавца это
значит, что операция прошла в~терминале, товар отдан, но через карточную сеть
деньги уже не получить, и~взыскивать придётся иначе.

Если же авторизации не существовало вовсе, эмитент всё равно списывает
по~полученной клиринговой записи без предварительного холда~--- сам клиринг
при~этом, как и~в~штатном сценарии, инициирован эквайером. Точкой списания
становится поступление клиринговой записи, а~не авторизация.
Риск переходит к~эквайеру и~продавцу; именно этот сценарий описывает force post
в~разделе~\ref{sec:tx-edge-cases}~\cite{mc-tpr}.

\subsection*{Расчёт interchange и~комиссий}

На этапе клиринга сеть рассчитывает interchange для каждой транзакции~---
межбанковскую комиссию, которую эквайер платит эмитенту;
см.~раздел~\ref{sec:ecosystem-interchange}. Ставка interchange зависит от
параметров, зафиксированных в~клиринговой записи: типа карты, MCC, способа
приёма, географии операции и ряда дополнительных признаков. Тип карты бывает
дебетовый, кредитный или корпоративный; способ приёма~--- чип, магнитная полоса
или CNP; география~--- внутренняя или международная.

Параллельно сеть начисляет собственные комиссии: плату за обработку, за
международную маршрутизацию, за использование специальных сервисов.
В~документах самих сетей эту составляющую называют scheme fees и~assessments. По
итогам обработки всех клиринговых записей за период сеть формирует чистые
расчётные позиции (net settlement positions) для каждого банка-участника~---
итоговые суммы, которые один банк должен другому с~учётом взаимозачёта
встречных потоков.

\section{Межбанковский расчёт (settlement)}\label{sec:tx-settlement}

\highlight{Расчёт завершает жизненный цикл транзакции~--- именно здесь деньги
переходят между банками}. До этого момента средства не покидали счёт держателя;
авторизация лишь заблокировала сумму, клиринг сформировал требование. На этапе
расчёта требование исполняется в~расчётной инфраструктуре платёжной
системы~\cite{fz-161,visa-core-rules,mc-tpr}.

\subsection*{Нетто-расчёт}

Карточные сети используют нетто-расчёт. Вместо отдельной межбанковской операции
на каждую транзакцию сеть суммирует все встречные обязательства между парами
банков за расчётный период и~переводит только разницу. Если за день банк А~как
эквайер должен банку~Б как эмитенту миллион рублей interchange, а~банк~Б как
эквайер должен банку~А как эмитенту восемьсот тысяч, фактически переводятся
только двести тысяч от~А~к~Б. Этот механизм многократно сокращает количество
и~объём межбанковских проводок.

Сеть проводит расчёт через назначенный расчётный банк или иной расчётный
сервис, предусмотренный её правилами. У Visa локальные операции, если в~стране
доступен соответствующий сервис, проходят через National Net Settlement Service
(NNSS), остальные~--- через International Settlement Service; правила отдельно
определяют расчётный банк и~расчётный счёт участника. У Mastercard итоговый
документ дня~--- ежедневное уведомление по~нетто-расчётам, и~оно считается
окончательной записью о~движении средств между
участниками~\cite{visa-core-rules,mc-tpr}.

\subsection*{Когда продавец получает деньги}

\emph{Расчёт} (settlement)~--- движение средств между банками внутри платёжной
системы; этим завершается жизнь транзакции в~сети. \emph{Выплата по~договору
эквайринга} (payout)~--- отдельный шаг: эквайер, уже получив средства от~сети,
перечисляет их продавцу по~условиям договора, ежедневно, через день или
по~иному графику, за~вычетом MSC. Четыре фазы и~разрывы между ними сведены
на~рис.~\ref{fig:tx-lifecycle-timeline}. В~продуктовом реестре эти два события
соответствуют разным переходам и~собственному графику выплат
(см.~раздел~\ref{sec:ledger-balances}).

Универсального срока для выплаты нет: после расчёта в~сети эквайер перечисляет
деньги продавцу по~графику договора~--- например, на~следующий рабочий день или
по~SLA (Service Level Agreement, договор об~уровне обслуживания) конкретного
провайдера. Для продавца с~большим оборотом или узкой маржой каждый
дополнительный день задержки замораживает оборотные средства.

\begin{figure}[tbp]
  \centering
  \includefiguresvg[width=\linewidth]{assets/figures/ch05-transaction-lifecycle/tx-lifecycle-timeline}
  \caption{Четыре фазы карточной транзакции и~три разрыва между ними.}\label{fig:tx-lifecycle-timeline}
\end{figure}

\section{Одностадийная и~двухстадийная схема}\label{sec:tx-dual-single}

\highlight{Описанная выше схема двухстадийна, \textit{dual message system}
(DMS): авторизация и~клиринг идут разными сообщениями в~разное время}. Сначала
терминал спрашивает \enquote{можно?}, потом эквайер предъявляет финансовую
запись. Это стандарт для большинства карточных
покупок~\cite{visa-core-rules,mc-tpr}.

Есть и~другая архитектура~--- одностадийная, \textit{single message system}
(SMS), в~которой авторизация и~финансовая запись совмещены в~одном
сообщении. Такой запрос одновременно спрашивает разрешение \emph{и} фиксирует
финансовое обязательство. Ответ эмитента означает сразу \enquote{одобрено}
и~\enquote{деньги можно считать списанными}, а~отдельного клирингового шага не
требуется.

\subsection*{Где используется каждая схема}

Двухстадийная схема доминирует в~обычных карточных покупках: в~магазинах,
ресторанах, во многих интернет-операциях. Она допускает расхождение между
авторизованной и~финальной суммой~--- это нужно там, где возможны чаевые,
довзвешивание товара или частичная отмена. При этом эквайер группирует транзакции
в~пакеты и~обрабатывает их сериями.

Одностадийная схема~--- стандартная модель для банкоматов (ATM). Когда
держатель снимает наличные, разделять авторизацию и~финансовую запись нет
смысла: деньги уже выданы физически, и~задним числом операцию не отменить. По
той же логике одностадийная схема используется в~некоторых дебетовых сетях, где
транзакция проходит на~точную сумму и~последующей
корректировки не предполагает~\cite{visa-core-rules,mc-tpr}.

На уровне ISO~8583 различие выражается в~MTI. Авторизационный запрос
в~двухстадийной схеме имеет MTI~0100, финансовый запрос в~одностадийной~---
MTI~0200; соответствующие ответы~--- MTI~0110 и~MTI~0210. Структура MTI
подробно разбирается в~разделе~\ref{sec:iso8583-mti}.

\importantnote{%
  Двухстадийная или одностадийная схема~--- свойство канала
  приёма, его не~выбирают ни~продавец, ни~разработчик.
  POS-терминал по~протоколу Visa или Mastercard отправляет
  двухстадийный поток; банкомат в~сети ATM~--- одностадийный.
}

\subsection*{Гибридные варианты}

Различие между двухстадийной и~одностадийной схемой размывается. Visa,
Mastercard и~процессоры поверх них поддерживают гибридные сценарии, в~которых
одобренная авторизация при определённых условиях автоматически доводится
до~финансовой записи или обслуживается общей инфраструктурой авторизации
и~клиринга. Гибриды снижают трудозатраты на стороне эквайера, но добавляют
сложности процессинговым системам, которым приходится поддерживать оба сценария
в~одной кодовой базе.

\section{Три способа отменить платёж}\label{sec:tx-cancellation}

Карточная транзакция в~каждый момент находится в~одной из фаз, и~способ отмены
определяется фазой. Этим карточная отмена отличается от банковского перевода,
который либо прошёл, либо нет.

\subsection*{Сообщение отмены (reversal)}

Отмена, reversal, аннулирует авторизацию \emph{до того}, как транзакция ушла
в~клиринг. В~терминах ISO~8583 это отдельный тип сообщения. MTI~0400, Reversal
Request, требует ответа от эмитента; MTI~0420, Reversal Advice, односторонне
уведомляет о~состоявшейся отмене. Эквайер отправляет одно из них в~сеть
с~просьбой аннулировать ранее одобренную авторизацию.

В~reversal-сообщениях DE~39 несёт код причины отмены, а~не код ответа;
конкретные значения~--- таймаут на стороне эмитента, ошибка разбора ответной
криптограммы~ARPC (Authorization Response Cryptogram),
обрыв связи и~прочие~--- задаются спецификацией сети или процессора.

Отмена возникает в~нескольких типовых сценариях. Покупатель передумал у~кассы;
терминал не получил ответ на авторизацию по таймауту и~отправляет отмену,
чтобы закрыть возможное одобрение; кассир ошибся с~суммой и~отменяет операцию.
При успешной отмене эмитент снимает холд со счёта держателя, и~средства возвращаются
в~доступный баланс обычно мгновенно или за несколько минут. Правила Mastercard
требуют от эквайера обеспечить отправку Reversal
Request/0400 в~течение 24~часов после отмены ранее авторизованной операции или
её завершения на меньшую сумму; у~Visa отмена также обязательна, если одобрение
пришло уже после таймаута~\cite{visa-core-rules,visa-tadg,mc-tpr}.

Отмена работает в~реальном времени, без финансовой записи и~без возвратной
транзакции: деньги ещё не~перемещались. Окно применимости закрывается с~уходом
транзакции в~клиринг: попавшую в~клиринговый файл авторизацию reversal уже не
отменит~\cite{mc-rules}.

\subsection*{Аннулирование до расчёта (void)}

\textit{void}, или аннулирование,~--- не отдельный тип сообщения карточной
сети, а~термин эквайринга и~кассового ПО для отмены операции до~окончательного
клиринга. В~\textit{dual-message}-среде отмена (\textit{cancellation}) должна
произойти до того, как эквайер отправит транзакцию в~клиринг; если отмена
случилась уже после выхода на онлайн-авторизацию, устройство или продавец
обязаны сформировать отмену, а~саму транзакцию убрать из пакета клиринга или
пометить как \textit{void}~\cite{visa-tadg}.

Для держателя карты \textit{void} и~отмена выглядят одинаково~--- холд
снимается, деньги возвращаются на счёт. На технической стороне \textit{void}
обрабатывается медленнее, потому что он привязан к~циклу клиринга, и~если файл
уже отправлен, корректировка попадёт только в~следующий батч.

\practicenote{%
  В~API большинства PSP и~платёжных шлюзов вы встретите метод
  \texttt{cancel} или \texttt{void}, абстрагирующий оба механизма.
  Если транзакция ещё не ушла в~клиринг, PSP отправляет отмену;
  если ушла, но расчёт не произошёл, применяется \textit{void}.
  Для разработчика, интегрирующего платежи, это удобная абстракция,
  но за ней стоят разные процессы с~разной скоростью
  возврата средств клиенту.
}

\subsection*{Возврат после расчёта (refund)}

Возврат~--- единственный способ вернуть деньги после того, как расчёт завершён
и~средства фактически перемещены между банками. В~отличие от отмены
и~\textit{void}, которые аннулируют исходную транзакцию, возврат~---
\emph{новая} транзакция в~обратном направлении: эквайер инициирует перевод от
продавца к~держателю через ту же карточную сеть.

Возврат проходит собственный цикл авторизации, клиринга и~расчёта, только
в~обратную сторону. В~Visa TADG (Transaction Acceptance Device Guide) возврат
описан как онлайн-авторизационное сообщение с~последующим клирингом; Mastercard
описывает возврат как операцию в \textit{dual-message}- или
\textit{single-message}-режиме, а~для возвратов в \textit{dual-message}-режиме
требует онлайн-авторизацию~\cite{visa-tadg,mc-tpr}.

Для держателя карты средства по возврату поступают на счёт через несколько
рабочих дней: возврат проходит полный цикл клиринга и~расчёта. Для продавца
возврат обычно дороже отмены: часть расходов исходной покупки автоматически не
компенсируется; эквайер нередко берёт отдельную комиссию за саму операцию
возврата. Все три механизма~--- отмена, \textit{void} и~возврат~--- расположены
относительно фаз транзакции на~рис.~\ref{fig:tx-cancellation-timeline}.

\begin{figure}[tbp]
  \centering
  \includefiguresvg[width=\linewidth]{assets/figures/ch05-transaction-lifecycle/tx-cancellation-timeline}
  \caption{Отмена, \textit{void} и~возврат на временной шкале транзакции.}\label{fig:tx-cancellation-timeline}
\end{figure}

\importantnote{%
  Распространённая ошибка при проектировании платёжной интеграции~---
  использовать возврат там, где достаточно отмены или \textit{void}.
  Если клиент отменяет заказ через минуту после оплаты, а~система
  продавца отправляет возврат вместо отмены, клиент будет ждать
  возврата несколько дней вместо нескольких минут, продавец потеряет
  interchange, а~эквайер зафиксирует лишнюю транзакцию. Практическое правило:
  отменять как можно раньше. Отмена лучше \textit{void}, \textit{void} лучше
  возврата.
}

\section{Граничные случаи}\label{sec:tx-edge-cases}

Схема \enquote{авторизация на~точную сумму, клиринг на~ту же сумму, расчёт
через день} покрывает основную массу розничных покупок. Остальные сценарии~---
предавторизация, инкрементальная и~частичная авторизация, отложенная
авторизация и~принудительная отправка~--- отличаются своими сроками
и~распределением риска.

\subsection*{Предавторизация (pre-authorization)}

Предавторизация, \textit{pre-auth},~--- авторизация, при которой финальная
сумма заранее неизвестна. Типовые примеры~--- гостиница и~аренда
автомобиля. В~момент заселения или получения машины продавец авторизует
оценочную сумму~--- ожидаемую стоимость проживания или депозит за
автомобиль,~--- а~итоговая определяется при выезде или возврате и~может
оказаться как меньше, так и~больше авторизованной.

В~сетевых правилах \textit{pre-auth}~--- самостоятельный тип запроса. Visa
различает Estimated Authorization и~Incremental Authorization, а Mastercard
рекомендует кодировать запрос как \textit{preauthorization}, если сумма
оценочная или если операция может вообще не завершиться. Эмитент ставит холд,
заранее зная, что финальный клиринг может отличаться от исходной оценки.
Допустимые пределы расхождения зависят от MCC, типа ТСП и~конкретного сценария,
поэтому точные пороги берут из актуальных правил
сети~\cite{visa-tadg,visa-core-rules,mc-tpr}.

\subsection*{Инкрементальная авторизация (incremental authorization)}

Инкрементальная авторизация увеличивает сумму ранее одобренной авторизации без
отмены исходной. Продавец отправляет дополнительный авторизационный запрос,
ссылающийся на исходный через код авторизации или специальный идентификатор,
и~эмитент увеличивает холд на запрошенную
сумму~\cite{visa-tadg,visa-core-rules,mc-tpr}.

Типичный сценарий~--- гостиница, где гость продлевает пребывание или заказывает
дополнительные услуги. Вместо повторного цикла \enquote{авторизовать~---
отменить} продавец наращивает сумму инкрементально. Поддержка такого потока
зависит от сети, региона, MCC и~настроек эмитента или эквайера, поэтому при
проектировании интеграции нужно закладывать запасной сценарий~--- например,
новую авторизацию на дополнительную сумму.

\subsection*{Частичная авторизация (partial authorization)}

Частичная авторизация~--- сценарий, в~котором эмитент одобряет транзакцию
только на часть запрошенной суммы. Держатель пытается оплатить покупку на
тысячу рублей, но на карте доступно только семьсот. Вместо полного отказа
с~кодом \texttt{51} эмитент одобряет операцию на семьсот рублей; продавец
предлагает клиенту доплатить оставшиеся триста наличными или другой
картой~\cite{visa-tadg,mc-tpr}.

Частичная авторизация требует поддержки на всех уровнях. Терминал обрабатывает
частичное одобрение и~предлагает доплату, POS-система корректирует чек, кассир
должен понимать, что происходит. Чаще всего схема встречается с~предоплаченными
и~подарочными картами, где остаток заранее известен и~ограничен.

\importantnote{%
  Если ваша система не поддерживает partial authorization, убедитесь,
  что терминал или платёжный шлюз не заявляет такую поддержку
  в~сетевом индикаторе partial approval. Например,
  у~Mastercard для этого используется Partial Approval
  Terminal Support Indicator в~DE~48, subelement~61.
  Получить частичное одобрение и~не обработать его корректно~---
  верный путь к~расхождениям в~учёте, когда продавец отпустит товар
  на полную сумму, а~спишется только часть~\cite{mc-tpr}.
}

\subsection*{Отложенная авторизация (deferred authorization)}

Отложенная авторизация~--- механизм, при котором транзакция совершается без
онлайн-авторизации, а~авторизационный запрос отправляется позже, по
восстановлении связи. Классические сценарии~--- оплата на борту самолёта, где
терминал работает офлайн и~авторизации уходят после посадки; паркинг, где сумма
определяется по времени стоянки; платные дороги и~общественный транспорт.

В~режиме deferred authorization терминал собирает данные карты, фиксирует
операцию локально и~позже формирует авторизационный запрос с~соответствующим
флагом. Эмитент может одобрить или отклонить этот запрос, но товар или услуга
уже предоставлены, и~кредитный риск несут эквайер и~продавец. Поэтому deferred
authorization разрешена только для определённых MCC и~при соблюдении лимитов на
сумму отдельной операции, устанавливаемых сетями. У Visa такие авторизации
нужно завершить в~течение 24~часов с~момента исходной операции, а~для
ряда наземных транспортных MCC~--- 4111,~4112,~4131~--- допускается
до~4~суток. В~авторизационном сообщении передаётся Deferred Authorization
Indicator со значением \texttt{5206} в~поле Field~63.3, обязателен для
card-present и~card-absent сценариев~\cite{visa-tadg,visa-core-rules}.

\subsection*{Принудительная отправка (force post)}

Принудительная отправка~--- механизм, при котором продавец отправляет
транзакцию в~клиринг без предварительной онлайн-авторизации, используя код
авторизации, полученный по~телефону через голосовую авторизацию (voice
authorization).

Такой сценарий возникает, когда онлайн-авторизация невозможна: терминал не
работает или пропала связь, но провести платёж нужно. Продавец получает голосовую
авторизацию у~эмитента или через эквайера и~затем использует этот код при
обработке транзакции. Mastercard оговаривает: код, полученный по каналу
голосовой авторизации, не защищает от \textit{chargeback} по
авторизационным основаниям~\cite{mc-tpr}.

Онлайн-проверки криптограммы здесь нет, часть контроля переносится на человека
и~процедуры эквайера, ответственность за последствия ложится на продавца. Force
post остаётся резервным сценарием: устойчивые каналы связи и~офлайн-режимы
терминалов закрывают большую часть случаев, где раньше требовалась голосовая
авторизация.

\begin{table}[htbp]
  \caption{Граничные случаи карточной авторизации.}\label{tab:tx-edge-cases}
  \begingroup
  \footnotesize
  \setlength{\tabcolsep}{4pt}
  \renewcommand{\arraystretch}{1.2}
  \begin{tabularx}{\textwidth}{
      @{}B{0.14\textwidth}
      P{0.17\textwidth}
      P{0.15\textwidth}
      P{0.17\textwidth}
      Y@{}
    }
    \toprule
    \headrow
    \thd{Механизм} & \thd{Когда сумма известна} &
    \thd{Онлайн-авторизация} & \thd{Кто несёт риск} &
    \thd{Типичный MCC} \\
    \midrule
    Pre-auth & позже & да
    & эквайер/ТСП & 7011 (гостиницы), 7512 (аренда авто) \\
    Incremental auth & постепенно & да (несколько)
    & эквайер/ТСП & 7011 \\
    Partial auth & сразу, но покрытие неполное & да
    & ТСП & любой (чаще prepaid) \\
    Deferred auth & позже & нет (офлайн)
    & эквайер/ТСП & 4511 (авиа), 7523 (паркинг) \\
    Force post & сразу & голосовая
    & ТСП & любой \\
    \bottomrule
  \end{tabularx}%
  \endgroup
\end{table}

\subsection*{Покупка с~выдачей наличных (POS-кешбэк)}\label{sec:tx-pos-cashback}

С~термином \enquote{cashback} в~обиходе путаница. Чаще всего так называют
маркетинговый бонус, который эмитент возвращает на счёт держателя по итогам
периода; физические наличные при этом не выдаются. POS-кешбэк устроен иначе:
продавец выдаёт держателю карты часть запрошенной суммы наличными вместе
с~розничной покупкой. У~Visa это \enquote{Cash at POS}
или \enquote{Visa Cashback}, у~Mastercard \enquote{Purchase with Cashback}
(PWCB); в~ПС~Мир~--- \enquote{Выдача наличных при совершении покупки}. На
прикладном уровне POS-терминала действует максимальный лимит выдачи наличных
за~операцию: у~НСПК~Мир~--- 5000~руб., у~Visa~--- 200~USD в~национальном
эквиваленте; Mastercard задаёт лимит в~местной валюте по~правилам конкретного
рынка.

POS-кешбэк допустим только для контактных и~бесконтактных чиповых карт; ввод по
магнитной полосе и~ручной ввод PAN запрещены. Кешбэк всегда привязан к~покупке.
Авторизация всегда онлайн~--- floor limit (порог суммы, ниже которого терминал
вправе одобрить операцию офлайн) для операции принудительно равен нулю
(раздел~\ref{sec:emv-risk-mgmt}), а~CVM (Cardholder
Verification Method, метод верификации держателя карты) обязателен. Набор
допустимых методов различается у~сетей: Mastercard требует только онлайн-PIN;
Visa допускает онлайн-PIN, офлайн-PIN и~CDCVM (Consumer Device CVM, верификация
  на устройстве пользователя~--- например, отпечаток пальца или Face~ID на
смартфоне) на мобильном устройстве; подпись и~No~CVM запрещены у~обеих. Обе
сети разрешают кешбэк только во внутренних операциях~--- международная выдача
наличных через POS не предусмотрена~\cite{visa-tadg,mc-tpr}.

На уровне EMV покупка с~выдачей наличных кодируется через два поля суммы. Теги
передаются в~формате TLV спецификации EMV~\cite{emvco-book3}. \texttt{0x9F02} (Authorized Amount) содержит общую сумму
операции~--- покупку плюс наличные, \texttt{0x9F03} (Other Amount)~--- отдельно
сумму наличных. Тип операции несёт тег \texttt{0x9C}; конкретные значения,
маркирующие наличие наличной части, задаются спецификациями сетей и~различаются
между ними. Сумма наличных ограничена прикладным лимитом терминала, не
описанным в~EMV-ядре, и~задаётся правилами сети для конкретного рынка.

В~ISO~8583 DE~4 несёт общую сумму, а~отдельная сумма наличных передаётся
в~DE~54 как дополнительная сумма; в~DE~55 дублируются EMV-теги \texttt{0x9F02},
\texttt{0x9F03} и~\texttt{0x9C}. Эмитент может одобрить операцию полностью или
отказать только в~наличной части, оставив покупку. На уровне ответа DE~39 для
этого предусмотрены отдельные коды \enquote{наличная часть не одобрена, покупка
проходит}; их конкретные значения различаются между Visa, Mastercard и~ПС~Мир,
и~эквайер обычно конвертирует фирменный код сети в~универсальный для своего
хоста ответ. Технический реверсал отменяет операцию целиком; обычная отмена или
возврат, инициированные кассиром по требованию держателя, допустимы только
в~отношении части покупки~\cite{visa-tadg,mc-tpr}.

Interchange по~POS-кешбэку устроен иначе, чем по~ATM-снятию. Тарифные сетки
Visa и~Mastercard относят покупку с~выдачей наличных к~розничной
debit-программе, а~не к~cash advance: у~Mastercard в~американской публичной
таблице ставок строка \enquote{Purchases, Purchases with Cash-back and Unique
Prepaid Rate} объединяет покупку и~покупку с~кешбэком в~единый
тариф~\cite{mc-merchant-rates-2025-2026}; у~Visa наличная часть в~Consumer
Debit Acceptance Fee исключается из объёма продаж и~не относится к~отдельной
cash-advance-ставке. ПС~Мир в~межхостовых правилах задаёт
свою ставку для розничной выдачи; конкретные значения~--- в~приложении~4
\enquote{Тарифы} к~Правилам ПС~Мир, которое публикует
НСПК~\cite{nspk-mir-rules}. Для ТСП кешбэк сокращает выручку наличными, которую
иначе пришлось бы инкассировать. В~Великобритании услугу поддерживает
государственная программа защиты доступа к~наличным: Закон~\enquote{Financial
Services and Markets Act~2023} наделил~FCA полномочиями обеспечивать разумное
наличие cash deposit и~withdrawal services, а~оператор крупнейшей банкоматной
сети~LINK ведёт реестр точек
cashback-without-purchase~\cite{uk-fsma-2023,fca-access-to-cash}. В~Канаде
POS-кешбэк действует с~1990-х как штатная функция дебетовой сети~Interac
и~принимается основными продуктовыми ритейлерами~\cite{interac-debit-cashback}.

Для держателя сумма наличных суммируется с~дневным лимитом снятия наличных
по~карте, если эмитент объединяет лимиты ATM и~POS; часть банков ведёт их
раздельно. Когда общего лимита не~хватает на~весь запрос, эмитент вправе
одобрить только покупочную часть и~отклонить наличную: правила Mastercard для
этого определяют отдельный код ответа DE~39 \texttt{87} \enquote{Purchase
Amount Only, No Cash Back Allowed}~\cite{mc-tpr}. Эмитент не~обязан принимать
такой исход: он может отказать в~операции целиком кодом~\texttt{51}
\enquote{insufficient funds} или~\texttt{61} \enquote{exceeds withdrawal amount
limit}; тогда продавцу остаётся предложить держателю повторить операцию без
кешбэка. Структура суммы, по~которой эмитент принимает это решение, передаётся
в~авторизационном запросе в~DE~4, DE~54 и~DE~55.

Для учёта это не бонусная программа, а~гибрид покупки и~выдачи наличных.
Ошибиться в~разделении сумм значит перепутать лимиты, interchange, реверсалы
и~последующую сверку.
